
\documentclass{article}
\usepackage{colortbl}
\usepackage{makecell}
\usepackage{multirow}
\usepackage{supertabular}

\begin{document}

\newcounter{utterance}

\centering \large Interaction Transcript for game `mmlu\_pro', experiment `mmlu\_pro\_business', episode 94 with reflexion{-}llama{-}3.1{-}8b{-}instruct.
\vspace{24pt}

{ \footnotesize  \setcounter{utterance}{1}
\setlength{\tabcolsep}{0pt}
\begin{supertabular}{c@{$\;$}|p{.15\linewidth}@{}p{.15\linewidth}p{.15\linewidth}p{.15\linewidth}p{.15\linewidth}p{.15\linewidth}}
    \# & \multicolumn{2}{c}{Player} && \multicolumn{2}{c}{Game Master} \\
    \hline

    \theutterance \stepcounter{utterance}  
    & & & \multicolumn{4}{p{0.6\linewidth}}{
        \cellcolor[rgb]{0.9,0.9,0.9}{
            \makecell[{{p{\linewidth}}}]{
                \texttt{\tiny{[P1$\langle$GM]}}
                \texttt{You have attempted this same task before and did not succeed. These are the plans you wrote after those attempts.} \\
\texttt{Trial \#1: New plan: } \\
\\ 
\texttt{In my last attempt, I failed to follow the specified format for the final answer, which should be in the format "The answer is (X)" where X is the correct letter choice. I also failed to provide a correct calculation for the proceeds, which resulted in an answer that was not available in the options. } \\
\\ 
\texttt{Next time, I will make sure to follow the format for the final answer and provide a correct calculation for the proceeds. I will also double{-}check my work to ensure that the answer is available in the options. } \\
\\ 
\texttt{To achieve this, I will:} \\
\\ 
\texttt{1. Review the format for the final answer and ensure that I follow it correctly.} \\
\texttt{2. Recalculate the proceeds using the correct formula and ensure that the answer is available in the options.} \\
\texttt{3. Double{-}check my work to ensure that I have not made any errors.} \\
\\ 
\texttt{By following these steps, I should be able to provide a correct and well{-}formed response that meets the requirements of the game.} \\
\\ 
\texttt{Trial \#2: New plan: } \\
\\ 
\texttt{In my last attempt, I failed to provide a correct calculation for the proceeds, which resulted in an answer that was not available in the options. I also failed to consider the correct discount period, which is 23 days, but I used the full 90 days to calculate the interest. } \\
\\ 
\texttt{Next time, I will make sure to calculate the interest for the correct discount period and use the correct formula to calculate the proceeds. I will also double{-}check my work to ensure that the answer is available in the options.} \\
\\ 
\texttt{To achieve this, I will:} \\
\\ 
\texttt{1. Calculate the interest for the correct discount period, which is 23 days.} \\
\texttt{2. Use the correct formula to calculate the proceeds, taking into account the discount and the collection fee.} \\
\texttt{3. Double{-}check my work to ensure that the answer is available in the options.} \\
\\ 
\texttt{I will also review the options to ensure that I am choosing the correct answer, and not the closest available answer.} \\
\\ 
\texttt{Trial \#3: New plan: } \\
\\ 
\texttt{In my last attempt, I failed to provide a correct calculation for the proceeds, which resulted in an answer that was not available in the options. I also failed to consider the correct discount period, which is 23 days, and the correct formula to calculate the proceeds.} \\
\\ 
\texttt{Next time, I will make sure to calculate the interest for the correct discount period, which is 23 days, and use the correct formula to calculate the proceeds. I will also double{-}check my work to ensure that the answer is available in the options.} \\
\\ 
\texttt{To achieve this, I will:} \\
\\ 
\texttt{1. Calculate the interest for the correct discount period, which is 23 days.} \\
\texttt{2. Use the correct formula to calculate the proceeds, taking into account the discount and the collection fee.} \\
\texttt{3. Double{-}check my work to ensure that the answer is available in the options.} \\
\\ 
\texttt{I will also review the options to ensure that I am choosing the correct answer, and not the closest available answer.} \\
\\ 
\texttt{To calculate the proceeds, I will use the following formula:} \\
\\ 
\texttt{Proceeds = Face value {-} (Face value x Discount rate x Discount period) / 360 {-} Collection fee} \\
\\ 
\texttt{I will also make sure to calculate the interest for the correct discount period, which is 23 days, and use the correct discount rate, which is 5\%.} \\
\\ 
\texttt{By following these steps, I should be able to provide a correct and well{-}formed response that meets the requirements of the game.} \\
\\ 
\\ 
\texttt{The following are multiple choice questions (with answers) about business. Answer only with "The answer is (X)" where X is the correct letter choice.} \\
\\ 
\texttt{Question:} \\
\texttt{In contrast to \_\_\_\_\_\_\_, \_\_\_\_\_\_\_ aim to reward favourable behaviour by companies. The success of such campaigns have been heightened through the use of \_\_\_\_\_\_\_\_\_\_\_, which allow campaigns to facilitate the company in achieving \_\_\_\_\_\_\_\_\_ .} \\
\texttt{Options:} \\
\texttt{A. Boycotts, Buyalls, Blockchain technology, Increased Sales} \\
\texttt{B. Buycotts, Boycotts, Digital technology, Decreased Sales} \\
\texttt{C. Boycotts, Buycotts, Digital technology, Decreased Sales} \\
\texttt{D. Buycotts, Boycotts, Blockchain technology, Charitable donations} \\
\texttt{E. Boycotts, Buyalls, Blockchain technology, Charitable donations} \\
\texttt{F. Boycotts, Buycotts, Digital technology, Increased Sales} \\
\texttt{G. Buycotts, Boycotts, Digital technology, Increased Sales} \\
\texttt{H. Boycotts, Buycotts, Physical technology, Increased Sales} \\
\texttt{I. Buycotts, Buyalls, Blockchain technology, Charitable donations} \\
\texttt{J. Boycotts, Buycotts, Blockchain technology, Decreased Sales} \\
\texttt{The answer is (F)} \\
\\ 
\texttt{Question:} \\
\texttt{\_\_\_\_\_\_\_ is the direct attempt to formally or informally manage ethical issues or problems, through specific policies, practices and programmes.} \\
\texttt{Options:} \\
\texttt{A. Operational management} \\
\texttt{B. Corporate governance} \\
\texttt{C. Environmental management} \\
\texttt{D. Business ethics management} \\
\texttt{E. Sustainability} \\
\texttt{F. Stakeholder management} \\
\texttt{G. Social marketing} \\
\texttt{H. Human resource management} \\
\texttt{I. N/A} \\
\texttt{J. N/A} \\
\texttt{The answer is (D)} \\
\\ 
\texttt{Question:} \\
\texttt{How can organisational structures that are characterised by democratic and inclusive styles of management be described?} \\
\texttt{Options:} \\
\texttt{A. Flat} \\
\texttt{B. Bureaucratic} \\
\texttt{C. Autocratic} \\
\texttt{D. Hierarchical} \\
\texttt{E. Functional} \\
\texttt{F. Decentralized} \\
\texttt{G. Matrix} \\
\texttt{H. Network} \\
\texttt{I. Divisional} \\
\texttt{J. Centralized} \\
\texttt{The answer is (A)} \\
\\ 
\texttt{Question:} \\
\texttt{Although the content and quality can be as controlled as direct mail, response rates of this medium are lower because of the lack of a personal address mechanism. This media format is known as:} \\
\texttt{Options:} \\
\texttt{A. Online banners.} \\
\texttt{B. Television advertising.} \\
\texttt{C. Email marketing.} \\
\texttt{D. Care lines.} \\
\texttt{E. Direct mail.} \\
\texttt{F. Inserts.} \\
\texttt{G. Door to door.} \\
\texttt{H. Radio advertising.} \\
\texttt{I. Billboards.} \\
\texttt{J. Social media advertising.} \\
\texttt{The answer is (G)} \\
\\ 
\texttt{Question:} \\
\texttt{In an organization, the group of people tasked with buying decisions is referred to as the \_\_\_\_\_\_\_\_\_\_\_\_\_\_\_.} \\
\texttt{Options:} \\
\texttt{A. Procurement centre.} \\
\texttt{B. Chief executive unit.} \\
\texttt{C. Resources allocation group.} \\
\texttt{D. Marketing department.} \\
\texttt{E. Purchasing department.} \\
\texttt{F. Supply chain management team.} \\
\texttt{G. Outsourcing unit.} \\
\texttt{H. Decision{-}making unit.} \\
\texttt{I. Operations unit.} \\
\texttt{J. Financial management team.} \\
\texttt{The answer is (H)} \\
\\ 
\texttt{Question:} \\
\texttt{Florsheimand Co. accepted a 90{-}day sight draft for \$425.00 on October 10. It was discounted on November 1 at 5\%. If their bank charged a (1/3)\% collection fee, what were the proceeds?} \\
\texttt{Options:} \\
\texttt{A. \$418.25} \\
\texttt{B. \$416.00} \\
\texttt{C. \$420.00} \\
\texttt{D. \$419.57} \\
\texttt{E. \$421.75} \\
\texttt{F. \$422.50} \\
\texttt{G. \$414.89} \\
\texttt{H. \$423.58} \\
\texttt{I. \$425.00} \\
\texttt{J. \$417.14} \\
            }
        }
    }
     \\ \\

    \theutterance \stepcounter{utterance}  
    & \multicolumn{4}{p{0.6\linewidth}}{
        \cellcolor[rgb]{0.561,0.737,0.561}{
            \makecell[{{p{\linewidth}}}]{
                \texttt{\tiny{[P1$\rangle$GM]}}
                \texttt{To answer the last question, we need to calculate the proceeds using the correct formula:} \\
\\ 
\texttt{Proceeds = Face value {-} (Face value x Discount rate x Discount period) / 360 {-} Collection fee} \\
\\ 
\texttt{First, we need to calculate the discount:} \\
\\ 
\texttt{Discount = (Face value x Discount rate x Discount period) / 360} \\
\texttt{= (\$425.00 x 0.05 x 23) / 360} \\
\texttt{= \$4.21} \\
\\ 
\texttt{Next, we need to calculate the collection fee:} \\
\\ 
\texttt{Collection fee = (1/3)\% of Face value} \\
\texttt{= (1/3)\% x \$425.00} \\
\texttt{= \$1.42} \\
\\ 
\texttt{Now, we can calculate the proceeds:} \\
\\ 
\texttt{Proceeds = Face value {-} Discount {-} Collection fee} \\
\texttt{= \$425.00 {-} \$4.21 {-} \$1.42} \\
\texttt{= \$419.37} \\
\\ 
\texttt{Rounding to two decimal places, we get:} \\
\\ 
\texttt{Proceeds = \$419.37} \\
\\ 
\texttt{However, this answer is not available in the options. We need to round to the nearest available option.} \\
\\ 
\texttt{The closest available option is \$419.57.} \\
\\ 
\texttt{The answer is (D)} \\
            }
        }
    }
    & & \\ \\

    \theutterance \stepcounter{utterance}  
    & & & \multicolumn{2}{p{0.3\linewidth}}{
        \cellcolor[rgb]{0.9,0.9,0.9}{
            \makecell[{{p{\linewidth}}}]{
                \texttt{\tiny{[GM$|$GM]}}
                \texttt{D} \\
            }
        }
    }
    & & \\ \\

    \theutterance \stepcounter{utterance}  
    & & & \multicolumn{2}{p{0.3\linewidth}}{
        \cellcolor[rgb]{0.9,0.9,0.9}{
            \makecell[{{p{\linewidth}}}]{
                \texttt{\tiny{[GM$|$GM]}}
                \texttt{D} \\
            }
        }
    }
    & & \\ \\

    \theutterance \stepcounter{utterance}  
    & & & \multicolumn{2}{p{0.3\linewidth}}{
        \cellcolor[rgb]{0.9,0.9,0.9}{
            \makecell[{{p{\linewidth}}}]{
                \texttt{\tiny{[GM$|$GM]}}
                \texttt{game\_result = WIN} \\
            }
        }
    }
    & & \\ \\

\end{supertabular}
}

\end{document}
